\section{Introduction}

Integrating refugees into host labor markets is a central policy challenge with implications for employment, welfare systems, and public support for immigration \citep{brell_labor_2020}. Policymakers rely on a mix of dispersal and welfare policies to manage these challenges. Dispersal policies assign refugees to specific regions, often limiting their initial choices, while welfare policies influence incentives to work or relocate. Yet integration outcomes also depend on local labor demand, which is beyond direct policy control. Understanding how refugees respond to these forces, and in particular how labor demand and welfare generosity interact, is crucial for designing effective integration strategies \citep{adda2022}.

This paper studies how local labor demand and welfare benefits at the time of receiving protection shape refugees’ labor market integration and location choices. We exploit the Austrian context, where asylum-seekers are exogenously dispersed across states with restricted labor market access during the asylum process. Upon receiving protection, they immediately gain full labor market rights and freedom to relocate. This institutional setting provides exogenous variation in both initial placement and exposure to local labor demand and benefit generosity, generating quasi-experimental conditions for causal identification.

Our findings can be summarized in four main points. First, higher local labor demand at the time of protection substantially increases employment, though effects dissipate after about 2.5 years. Second, more generous welfare benefits slightly reduce employment, with similar fading over time. Third, both factors have persistent effects on internal migration: refugees are more likely to remain in their assigned location when labor demand is strong or benefits are high, and many otherwise relocate to Vienna. Finally, we document strong interactions: the employment effects of welfare generosity are amplified in regions with high labor demand. Together, these results highlight the short-run importance of local conditions for employment and the long-run role of policies in shaping refugee mobility.

Our paper contributes to several strands of literature. We add to work on immigrants’ location choices and the welfare magnet hypothesis \citep{borjas1999, ferwerda2023, monras2023}, showing that while welfare generosity matters, its effects are modest compared to labor demand. We also contribute to research on the role of initial conditions in shaping refugees’ integration trajectories \citep{aslund_when_2007, fasani_struggle_2022, dustmann2024}, demonstrating that initial labor demand shocks fade relatively quickly for employment but persist for mobility. By examining both employment and relocation jointly, we bridge two literatures often studied separately and provide new evidence on the interaction of local labor market conditions and welfare policy.

We use administrative data covering over 40,000 refugees arriving in Austria between 2011 and 2018, including detailed labor market trajectories, welfare receipt, and residential histories. Our empirical strategy exploits exogenous initial placement, restrictions during the asylum process, and rich spatial-temporal variation in both labor demand and benefit levels. This fine-grained data allows us to identify both independent and interactive effects of labor demand and welfare generosity on refugees’ integration choices. The remainder of the paper is structured as follows: Section \ref{sec:background} provides institutional background, Section \ref{sec:dataoob} describes the data, Section \ref{sec:empirics} outlines our empirical strategy, Section \ref{sec:results} presents the results, and Section \ref{sec:conclusion} concludes.